\documentclass[10pt,a4paper]{amsart}
\usepackage[margin=19mm]{geometry}
\usepackage{amsmath,amssymb,mathtools}
\usepackage[scr=rsfs,cal=euler]{mathalfa}
\usepackage{xfrac}
\usepackage[unicode,hidelinks,bookmarks=false]{hyperref}
\newcommand{\ie}{i.e.\ }
\newcommand{\cf}{cf.\ }
\newcommand{\resp}{respectively\ }
\newcommand{\Subsection}{\subsection}
\providecommand{\nobreakdash}{\nobreak\hbox{-}}
\setlength{\emergencystretch}{2em}
\multlinegap0pt
\usepackage{tikz}
\usepackage{mathtools}
\usepackage{bm}
\newtheorem{proposition}{Proposition}
\newcommand{\mirror}{\textsf{mirror}}
\title{Average signature and 4-genus of 2-bridge knots}
\author{Moshe Cohen, Adam M. Lowrance, Neal Madras, Steven Raanes}
\date{Source: arXiv:2406.17738v2 (CC BY 4.0)}
\begin{document}
\maketitle

\noindent\emph{Source and licence.} Average signature and 4-genus of 2-bridge knots. Version arXiv:2406.17738v2 (CC BY 4.0), distributed by arXiv under the Creative Commons Attribution 4.0 International licence, http://creativecommons.org/licenses/by/4.0/, as recorded in the arXiv licence metadata for this exact version and shown on the abstract page https://arxiv.org/abs/2406.17738v2. Full licence notice: \texttt{licenses/arxiv-2406.17738-CC-BY-4.0.txt}.\par\smallskip

\noindent\textit{Editorial scope.} Counted unit: the article's proposition prop.Jbound (source0.tex lines 2971--2975). The article's complete original proof of it is delivered with it (source0.tex lines 2976--3068, one \texttt{proof} environment of 93 lines). No mathematics is written, repaired or re-derived: the statement and the proof are the article's own lines, in the article's own notation, and no retained line is substituted or re-flowed.  Retained uncounted as the source-local context the proof needs: lines 2918--2922.  Nothing was omitted from any retained span, so the package carries no omission record.  No retained line contains a citation, so the wrapper carries no bibliography and no bibliography entry.  No retained line contains a figure environment, an image-inclusion command or drawing code, so no image is delivered and the package declares no asset dependency.  Editorial formatting only, in addition: an AMS \texttt{amsart} preamble that restates the article's own declarations and macros used by the retained lines, namely the environments \texttt{proposition} on separate counters, together with the standard packages those lines need.  The retained environments print in this document as Proposition 1 in order of appearance, whereas the article prints the same items with the numbering of its own full document.  The complete native source of the article as distributed in the arXiv e-print for this exact version is bundled unchanged as \texttt{sources/161371/source0.tex}.
\begin{equation}
   \label{eq.pkdiffbound}
   \left| (\mathbf{P}^{k})(o_i,o_j)  \,-\,  (\mathbf{P}^{k})(o_i,o_\ell) \right|   \;\leq \;(2^{-s})^k
   \hspace{5mm}\hbox{for all }o_i,o_j,o_\ell \in OR.
\end{equation}

\begin{proposition}
   \label{prop.Jbound}
Let $s$ and $t$ be positive integers, and let $\vec{\omega}\in OW^{Chi}(s)$. The expected absolute value of the displacement from the origin in the $\vec{\omega}$ coordinate of the word-counting process after $t$ steps satisfies 
$E| D_{\vec{\omega}}(t)|   \,\leq \,  2\sqrt{t/2^s}$.   
\end{proposition}

\begin{proof}
For our fixed $\vec{\omega}\in OW^{Chi}(s)$, we shall write
 $\vec{\omega}$ as $(o_*,z_*)$ and 
 $\mirror(\vec{\omega})$ as $(\overline{o}_*,\overline{z}_*)$.  
Recall that we write $o_*'$ for the orientation that occurs immediately after $\vec{\omega}$, and then $\overline{o}_*'$ is the orientation that occurs immediately after $\mirror(\vec{\omega})$.

Since $\left(E|D_{\vec{\omega}}(t)|\right)^2 \;\leq \;  E(D_{\vec{\omega}}(t)^2)$ by the Cauchy-Schwarz inequality, our goal shall be to obtain a bound on $E(D_{\vec{\omega}}(t)^2)$.  We use 
\begin{equation}
   \label{eq.JsumI}
     E(D_{\vec{\omega}}(t)^2)  \;=\;  \sum_{\ell=1}^t E(I_{\vec{\omega}}(\ell)^2)  \,+\,  2 \sum_{k=1}^{t-1} \sum_{\ell=1}^{t-k}  E( \,I_{\vec{\omega}}(\ell)\,I_{\vec{\omega}}(\ell+k)\, ) \,.
\end{equation}
We first consider $E(I_{\vec{\omega}}(\ell)^2)$.  Observe that $I_{\vec{\omega}}(\ell)$ equals 0 unless $(X_\ell,Z_\ell)$ is either $(o_*,z_*)$ or $(\overline{o}_*,\overline{z}_*)$, in which 
case $I_{\vec{\omega}}(\ell)^2$ equals 1.
Thus
\begin{eqnarray} 
    \nonumber
    E(I_{\vec{\omega}}(\ell)^2)  &=&  \Pr(X_\ell=o_*,Z_\ell=z_*)+\Pr(X_\ell=\overline{o}_*,Z_\ell=\overline{z}_*)   \\
    \label{eq.I2bound}
  & \leq &    \Pr(Z_\ell=z_*)+\Pr(Z_\ell=\overline{z}_*) \qquad\;=\;  \frac{2}{2^{s}}.
\end{eqnarray}


Next  we consider the term $ E( \,I_{\vec{\omega}}(\ell)\,I_{\vec{\omega}}(\ell+k)\, )$, where $k\geq 1$.
Observe that $I_{\vec{\omega}}(\ell)I_{\vec{\omega}}(\ell+k)$ equals 0 unless $(X_\ell,Z_\ell)$ and $(X_{\ell+k},Z_{\ell+k})$ each is either $(o_*,z_*)$ or $(\overline{o}_*,\overline{z}_*)$, in which 
case $I_{\vec{\omega}}(\ell)I_{\vec{\omega}}(\ell+k)$ equals $1$ or $-1$.
Thus 
\begin{equation}
    \label{eq.Eaaaa}
     E(\,I_{\vec{\omega}}(\ell)\,I_{\vec{\omega}}(\ell+k)\, )   \;=\;  a(1,1) \,-\,a(1,-1) \,+\,a(-1,-1)\,-\,a(-1,1) \,,
\end{equation}
where
\[    a(i,j)   \;=\;  \Pr(I_{\vec{\omega}}(\ell)=i, \,I_{\vec{\omega}}(\ell+k)=j)\,.
\]

We will need to exploit some cancellation here.  Let's start with $a(1,1)-a(1,-1)$.  We have
\[   a(1,1)  = \Pr(X_\ell=o_*,Z_\ell=z_*,X_{\ell+k}=o_*,Z_{\ell+k}=z_*)  
        = u_1 u_2 u_3 u_4\,,   
\]
where
\begin{eqnarray*}
u_1\; & =\; & \Pr(X_\ell=o_*),   \\
u_2 \; & =\; &  \Pr(Z_{\ell}=z_*\,|  \,  X_\ell=o_*)  \\  u_3  \; & =\; &  \Pr(X_{\ell+k}=o_*\,|  \,  X_\ell=o_*,Z_\ell=z_*)   
        \\   
u_4 \; & =\; & \Pr(Z_{\ell+k}=z_*\,|  \,  X_\ell=o_*,Z_\ell=z_*,X_{\ell+k}=o_*)   \,.
\end{eqnarray*}
By the independence of terms of the $Z$ sequence, we have $u_2=u_4 = 2^{-s}$.  
The value of $u_1$ depends on the initial orientation $X_1=o_1$, via
\[   u_1 \;=\;  \Pr(X_\ell=o_*)  \;=\;  
  \mathbf{P}^{\ell-1}(o_1,o_*) \,.
\]
The quantity $u_3$ is a bit more subtle.   Given the information that $X_\ell=o_*$ and $Z_\ell=z_*$, it must also be 
true that $X_{\ell+1}=o_*'$.  Therefore, the Markov property of the $X$ process tells us that 
\[  u_3 \;=\;  \Pr(X_{\ell+k}=o_*\,|  \,  X_{\ell+1}=o_*') \;=\;  (\mathbf{P}^{k-1})(o_*',o_*) \,.
\]
Putting the pieces together shows that 
\[     a(1,1)  \;=\;  (2^{-s})^2 \,\Pr(X_\ell=o_*)  \,(\mathbf{P}^{k-1})(o_*',o_*) \,.
\]
An exactly analogous argument (with the only difference being in the $u_3$ term) shows that 
\begin{eqnarray*}
    a(1,-1)  & = & \Pr(X_\ell=o_*,Z_\ell=z_*,X_{\ell+k}=\overline{o}_*,Z_{\ell+k}=\overline{z}_*)  \\
       & = & (2^{-s})^2 \,\Pr(X_\ell=o_*)  \,(\mathbf{P}^{k-1})(o_*',\overline{o}_*) \,.
\end{eqnarray*}
Then 
\begin{eqnarray}
   \nonumber
   |a(1,1)-a(1,-1)|  &=  & (2^{-s})^2 \,\Pr(X_\ell=o_*)  \,\left| (\mathbf{P}^{k-1})(o_*',o_*) \,-\,   (\mathbf{P}^{k-1})(o_*',\overline{o}_*)\right|
     \\
     \label{eq.adiff1}
     & \leq &  (2^{-s})^2 \cdot 1 \cdot (2^{-s})^{k-1}        \hspace{12mm}\hbox{(by Eq.\ (\ref{eq.pkdiffbound}))}.
\end{eqnarray}
Again, a very similar calculation shows that
\begin{equation}
           \label{eq.adiff2}
    |a(-1,-1)\,-\,a(-1,1)|  \;\leq \;  (2^{-s})^2 \cdot 1 \cdot (2^{-s})^{k-1}.
\end{equation}
Inserting Equations (\ref{eq.adiff1}) and (\ref{eq.adiff2}) into (\ref{eq.Eaaaa}) shows that 
\begin{equation}
   \label{eq.EIIbound}
      \left|   E(\,I_{\vec{\omega}}(\ell)\,I_{\vec{\omega}}(\ell+k)\, ) \right|  \;\leq \;  2 \,(2^{-s})^{k+1} \hspace{5mm} \hbox{for all $k\geq 1$}.
\end{equation}
Finally, using Equations (\ref{eq.JsumI}), (\ref{eq.I2bound}),     and (\ref{eq.EIIbound}), we obtain
\begin{eqnarray}
  \nonumber
   E(D_{\vec{\omega}}(t)^2) & \leq &   2t \,2^{-s} \,+\,  2\sum_{k=1}^{t-1}  (t-k)  \,(2^{-s})^{k+1}
   \\
    \nonumber
   & \leq & 2t  \left(2^{-s}  \,+  \,\frac{(2^{-s})^2}{1-2^{-s}} \right)
   \\
      \label{eq.J2sumbound}
   &\leq & 4t\,2^{-s}   \hspace{10mm}\hbox{(since $2^{-s}\leq 1/2$)}.
\end{eqnarray}
The proposition now follows because $E|D_{\vec{\omega}}(t)| \;\leq \;  \sqrt{E(D_{\vec{\omega}}(t)^2)}$.
\end{proof}

\end{document}
